\subsection{Rank Theorem}

\begin{perspective} \label{per:rank-theorem-linear-model} \hlt{Rank Theorem Linear Model}\\
Theorem \ref{thm:quotient-by-kernel-is-range} reduces a linear map to its kernel--range split.\\
The constant-rank theorem is the local $C^1$ analogue after coordinate changes.
\end{perspective}

\begin{theorem} \label{thm:constant-rank-theorem-rudin} \hlt{Rank Theorem}\\
Let $E\subseteq\R^n$ be open, let $\mathbf{F}\in C^1(E,\R^m)$ have constant rank $r$, and fix $\a\in E$. Set $A\defeq D\mathbf{F}(\a)$, $Y\defeq\operatorname{range}A$, and let $P:\R^m\to Y$ be the orthogonal projection. Then for some $\x_0\in\R^n$, neighborhoods $V_0\ni\x_0$ and $U\ni\a$, and $C^1$ diffeomorphism $H:V_0\to U$, there exists $\varphi\in C^1(A(V_0),\ker P)$ such that $\mathbf{F}\circ H=A+\varphi\circ A$ on $V_0$.
\begin{proof}
If $r=0$, Theorem \ref{thm:bounded-derivative-convex-mean-value} makes $\mathbf{F}$ locally constant; take $H=I$ and $\varphi(0)=\mathbf{F}(\a)$.\\
Assume $r>0$. Choose $S:Y\to\R^n$ with $AS=I_Y$ and set $G\defeq I+SP(\mathbf{F}-A)$. Since $DG(\a)=I$, Theorem \ref{thm:inverse-function-theorem-rn} gives a local $C^1$ diffeomorphism $G:U\to V_0$ with inverse $H$; shrink $V_0$ convex and put $\x_0\defeq G(\a)$.\\
Let $\Psi\defeq\mathbf{F}\circ H$. From $AG=P\mathbf{F}$, $P\Psi=A$ and $PD\Psi=A$.\\
Since $\operatorname{rank}D\Psi=\operatorname{rank}A=\operatorname{rank}PD\Psi=r$, $P$ is injective on $\operatorname{range}D\Psi$; hence $A\mathbf{v}=0$ implies $D\Psi(\x)\mathbf{v}=0$.\\
By Theorem \ref{thm:bounded-derivative-convex-mean-value}, $\Psi$ is constant on each convex fiber $V_0\cap(\x+\ker A)$.\\
Define $\varphi(A\x)\defeq(I-P)\Psi(\x)$. Fiber constancy makes $\varphi$ well-defined, $\varphi(A\x)\in\ker P$, and $\Psi=A+\varphi\circ A$. After shrinking $V_0$ so $\sigma(\mathbf{u})\defeq\x_0+S(\mathbf{u}-A\x_0)$ lies in $V_0$ for $\mathbf{u}\in A(V_0)$, $\varphi=(I-P)\Psi\circ\sigma$ is $C^1$.\\
\hlt{Structural Commentary}: Constant rank gives projection-plus-graph normal form.\\
\hlt{Key Insight}: Constant rank kills variation along $A$-fibers.
\end{proof}
\end{theorem}

Rank-nullity (Theorem \ref{thm:rank-nullity}) says that the domain directions split into rank directions and fiber directions. The rank theorem says a constant-rank map has, locally, the same split after $C^1$ changes of coordinates.

\begin{figure}[H]
  \centering
  \begin{tikzpicture}[
  scale=0.9,
  >=Stealth,
  every node/.style={font=\scriptsize},
  fiber/.style={blue!70!black, line width=0.75pt},
  image/.style={black!70, line width=0.75pt},
  arrow/.style={->, line width=0.75pt},
  point/.style={circle, fill=black, inner sep=1.1pt}
]
  \begin{scope}
    \node[anchor=south] at (1.55,2.85) {original coordinates};
    \foreach \x/\s in {0.45/0.05,0.95/-0.04,1.45/0.08,1.95/-0.02,2.45/0.04} {
      \draw[fiber]
        (\x,0.20) .. controls (\x-0.16,0.85+\s) and (\x+0.18,1.35-\s) .. (\x,2.55);
    }
    \node[left=2pt] at (0.35,1.38) {fibers};
  \end{scope}

  \draw[arrow] (3.10,1.38) -- node[above=2pt] {$\mathbf{F}$} (4.18,1.38);

  \begin{scope}[xshift=4.45cm]
    \node[anchor=south] at (1.45,2.85) {image};
    \draw[image] (0.10,1.00) .. controls (0.90,1.50) and (1.85,1.35) .. (2.90,2.15);
    \draw[image] (0.20,1.30) .. controls (1.05,1.82) and (1.95,1.65) .. (3.02,2.48);
    \foreach \x/\y in {0.42/1.13,0.92/1.31,1.42/1.48,1.92/1.66,2.42/1.89} {
      \node[point] at (\x,\y) {};
    }
    \node[right=2pt] at (2.88,2.36) {$\operatorname{range}D\mathbf{F}$};
  \end{scope}

  \begin{scope}[yshift=-3.35cm]
    \node[anchor=south] at (1.55,2.85) {normal form};
    \foreach \x in {0.45,0.95,1.45,1.95,2.45} {
      \draw[fiber] (\x,0.20) -- (\x,2.55);
    }
    \node[left=2pt] at (0.35,1.38) {$\ker A$-fibers};
  \end{scope}

  \draw[arrow] (3.10,-1.97) -- node[above=2pt] {$P$} (4.18,-1.97);

  \begin{scope}[xshift=4.45cm,yshift=-3.35cm]
    \node[anchor=south] at (1.45,2.85) {projection image};
    \draw[image] (0.25,1.38) -- (3.00,1.38);
    \foreach \x in {0.45,0.95,1.45,1.95,2.45} {
      \node[point] at (\x,1.38) {};
    }
    \node[right=2pt] at (3.00,1.38) {$\R^r$};
  \end{scope}
\end{tikzpicture}

  \caption{After local coordinate changes, a constant-rank map looks like projection onto the rank directions.}
  \label{fig:rank-theorem-projection-normal-form}
\end{figure}

\begin{theorem} \label{thm:lagrange-multiplier-one-constraint} \hlt{One-Constraint Lagrange Multiplier}\\
Let $E\subseteq\R^n$ be open, $f,g:E\to\R$ be $C^1$, and $S=\{\x:g(\x)=c\}$.\\
If $p\in S$, $\nabla g(p)\neq\ze$, and $f|_S$ has a local extremum at $p$, then $\exists\lambda\in\R$ with $\nabla f(p)=\lambda\nabla g(p)$.
\begin{proof}
If $Df(p)=0$, take $\lambda=0$. Otherwise suppose $Df(p)$ is not a multiple of $Dg(p)$.\\
Complete $Dg(p),Df(p)$ to a dual basis $Dg(p),\quad Df(p),\quad \ell_3,\ldots,\ell_n$ by Theorem \ref{thm:dual-basis-and-dimension},\\
and set $\Phi(\x)=\bigl(g(\x)-c,\ f(\x)-f(p),\ \ell_3(\x-p),\ldots,\ell_n(\x-p)\bigr)$.\\
Then $D\Phi(p)$ is invertible, so Theorem \ref{thm:inverse-function-theorem-rn} makes $\Phi$ locally invertible. In these coordinates, $S$ is locally the slice with first coordinate $0$, while the second coordinate is $f-f(p)$. Local invertibility gives nearby points of $S$ with this second coordinate of both signs, contradicting extremality.\\
Hence $Df(p)=\lambda Dg(p)$; Riesz representation, Theorem \ref{thm:riesz-representation-finite-dimensional}, gives the gradient form.\\
\hlt{Structural Commentary}: The constraint kills first-order variation along $S$.\\
\hlt{Key Insight}: If $Df$ were independent of $Dg$, $f$ would be a free constraint coordinate.
\end{proof}
\end{theorem}
